% Preambula zajednicka za glavni dokument i poglavlja.
\usepackage[T1]{fontenc}
\usepackage[utf8]{inputenc}
\usepackage[serbian]{babel}
\usepackage{lmodern}

\usepackage[a4paper,margin=3cm]{geometry}
\usepackage[patch=none]{microtype}
\usepackage{graphicx}
\usepackage{float}
\usepackage{amsmath,amssymb,amsthm}
\usepackage{mathtools}
\usepackage{enumitem}
\usepackage{hyperref}
\usepackage{fancyhdr}
\usepackage{setspace}

\hypersetup{
	colorlinks=true,
	linkcolor=black,
	citecolor=black,
	urlcolor=blue
}

\setlength{\parindent}{1.2em}
\setlength{\parskip}{0.4em}
\onehalfspacing{}
\setlength{\headheight}{14.5pt}
\addtolength{\topmargin}{-0.5pt}

% Fancy headers/footers with chapter titles (oneside).
\pagestyle{fancy}
\fancyhf{}
\makeatletter
\@ifundefined{chapter}{
	\renewcommand{\sectionmark}[1]{\markboth{Glava #1}{}}
}{
	\renewcommand{\chaptermark}[1]{\markboth{Glava \thechapter. #1}{}}
	\renewcommand{\sectionmark}[1]{\markright{#1}}
}
\makeatother
\fancyhead[L]{\MakeUppercase{\leftmark}}
\fancyhead[R]{\thepage}
\renewcommand{\headrulewidth}{0.4pt}
\renewcommand{\footrulewidth}{0pt}

% Use the same header style on chapter-opening pages.
\fancypagestyle{plain}{
	\fancyhf{}
	\fancyhead[L]{\MakeUppercase{\leftmark}}
	\fancyhead[R]{\thepage}
	\renewcommand{\headrulewidth}{0.4pt}
	\renewcommand{\footrulewidth}{0pt}
}

\graphicspath{{images/}}

% Okruzenja za teoreme, definicije i primere.
\newtheorem{theorem}{Teorema}[section]
\newtheorem{lemma}[theorem]{Lema}
\newtheorem{proposition}[theorem]{Propozicija}
\newtheorem{corollary}[theorem]{Posledica}

\theoremstyle{definition}
\newtheorem{definition}[theorem]{Definicija}
\newtheorem{example}[theorem]{Primer}

\theoremstyle{remark}
\newtheorem{remark}[theorem]{Napomena}


% Nove komande 
\newcommand{\ket}[1]{| #1 \rangle}
\newcommand{\latt}{\mathcal{L}}